\section{Introduction}
\label{sec:intro}

Questions about a recorded lecture usually come down to a single moment: the slide with the definition, the plot that compares two methods, the line of a derivation.
A vision-language model (VLM) can answer only if that moment is among the frames it is given.
Models with long contexts can take in hundreds of frames.
A model that fits on one 16\,GB consumer GPU cannot, and with a budget of a few frames the choice of frames becomes the problem.

Lectures come with an unusual hint.
Speakers narrate what they show (``this plot compares the two methods'', ``on the next slide we have the update rule''), so the timed transcript points, imperfectly, to the moments where an answer is visible.
Recent systems already use subtitles to guide frame retrieval~\cite{he2025vsi,wei2026gcr}, and learned frame selectors are now common~\cite{yu2025framevoyager,hu2025mllm,buch2025flexible}.
Three questions remain open.
How much does the speech buy when only a handful of frames can be shown?
Does a small trained component improve on carefully built rules?
And can answer accuracy, the measure almost everyone reports, tell such systems apart?

We study these questions under tight, fixed conditions.
A frozen answer model, usually Qwen3-VL-4B~\cite{bai2025qwen3vl}, receives at most four frames from the system under test, always with the same prompt.
The systems range from a starting design that searches transcript windows, through a tuned rule-based pipeline with a dense picture index and a diversity-aware picker, to two small trained heads: a speech-to-moment locator and a frame-set selector.
We also include training-free baselines, among them a rule that fuses subtitle and visual scores in the spirit of VSI~\cite{he2025vsi}.
Every method was fixed on development data before its test set was run, and the comparisons we test were declared in advance and corrected for multiplicity.
We evaluate on 894 held-out LongVideoBench questions~\cite{wu2024longvideobench} and on a lecture test set built from M3AV~\cite{chen2024m3av}.
In the lecture set, each question is generated from a slide and carries an automatically labelled interval in which that slide is on screen, so we can measure whether a system showed the right slide, not only whether the answer came out right.

\paragraph{Findings.}
\begin{enumerate}
\item \textbf{The pipeline helps as a whole; the trained heads add nothing measurable to it.}
On fresh LongVideoBench questions, the tuned pipeline with selected speech text beats the starting design by 8.9 points with Qwen3-VL-4B and by 8.3 points with a 120B answer model.
Adding the two trained heads gives no confirmed gain with either model, and the tuned pipeline itself shows no advantage over simply taking the four frames most similar to the question in a dense index.
\item \textbf{In lectures, speech-guided localization works, but a training-free rule keeps up with it.}
The trained locator shows the source slide for 79.4\% of test questions, 5.3 points more than the rule-based pipeline.
The VSI-inspired rule reaches 79.8\%; the difference of $-0.5$ points has a 95\% interval of $-3.8$ to $+2.9$, so neither is shown to be better.
\item \textbf{Accuracy agrees with evidence placement but sits on a high floor.}
The question alone answers 51.2\% of the lecture questions, and the transcript alone 62.1\%.
The source slide is lost about as often while gathering candidates as when picking the final four, and answers are right 93.9\% of the time when it is shown, against 56.1\% to 66.0\% when it is not.
\end{enumerate}

\paragraph{Contributions.}
We do not propose a new state-of-the-art selector.
Our contributions are empirical:
(1) a controlled evaluation of transcript-guided frame selection under a four-frame budget, with methods fixed before testing, held-out test sets, declared contrasts, no-video controls and a human check of a question sample;
(2) a comparison of rule-based, learned and training-free fusion selectors under one frozen answer model and one frame allowance, including negative results for two learned heads on fresh data;
(3) a stage-wise analysis that separates where evidence is lost from where answers go wrong, and the reporting practices it suggests.
\Cref{sec:related} places the study in context, \cref{sec:study} describes the systems, data and protocol, \cref{sec:results} presents the results, and \cref{sec:discussion} discusses what they mean and where they stop.
